\section{Introduction}
\vspace{0.5cm}
\begin{center}
    \textit{"A good portfolio is more than a long list of \\
    good stocks and bonds. It is a balanced whole $[...]$."}\\
    -- \textbf{H. Markowitz \cite{MarkowitzQuote}} --
\end{center}
\vspace{0.5cm}
    
\subsection{Motivation}
This document was made to be an introduction to quantitative finance through development of market risk analysis algorithms.
It focuses on the study of the binomial model, known as the Cox-Ross-Rubinstein model.

Firstly, we will talk about the theoretical foundations of this prediction model before analyzing how it differs from the Black-Scholes and Monte Carlo models.
Finally, we will implement and analyze the  market variation estimators, known as the "Greeks".

I am working on this paper in my spare time for personal learning and curiosity. Consequently, it may contain errors or inaccuracies. 
The provided sources are in French or English, you may use a translator if you do not understand either of these languages.

\subsection{History of option valuation}

Modern portfolio theory in finance originated in 1952 with the work published by Harry Markowitz \cite{HarryMarkowitz}.
This theory explains how investors can optimize their portfolios and analyze asset pricing based on market conditions (volatility, risk and so on). 
Therefore it lays the foundations of quantitative finance by defining the portfolio as a collection of financial assets.

Based on Markowitz's work, three economists \textit{Fischer Black}, \textit{Myron Scholes} and \textit{Robert Merton} published an article in 1973 presenting a formula allowing to calculate the theoretical price of a European option.
Their contribution was to find a partial differential equation (PDE) to find the fair price for a European call option. \cite{Option}

With the development of computers and the exponential growth of computing power, a new method emerged in 1977 : the Monte Carlo approach. 
This model involves randomly simulating thousands of possible price paths.
The option price is the the average of the payoffs obtained from these paths. \cite{MonteCarloPricing}
It is the primary method for exotic options such as Asian options, where the payoff depends on the average price over a certain period of time.\cite{AsianOption}

Two years later, a new trio of economists :  \textit{John Cox}, \textit{Stephen Ross} and \textit{Mark Rubinstein} proposed an innovative alternative : the binomial model, also known as the Cox-Ross-Rubinstein model.
In this model, time is discretized, allowing the path of an asset price to be represented as a binary tree.
This simplifies estimation, providing a numerical approximation to the Black-Scholes model which relies on an anlytical solution.
The approach shifts the problem from a complex partial differential equation to a simple approximation that can be programmed on a computer.
Furthermore, the model enables the valuation of American options, wich can be exercised at any time with only a minor change in the algorithm. \cite{OptionPricingASimplifiedApproach}

Subsequently, other more exotic models appeared (e.g. trinomial trees) \cite{TrinomialTree}.
In one of the following part, we will focus on the binomial model and compare it with the Black-Scholes and Monte Carlo models, representing respectively an approximation model, an analytical model and a simulation-based one.
The Cox-Ross-Rubinstein model is more intuitive and simpler to implement algorithmically than the Black-Scholes model.

\subsection{Goals}

The main goal of this paper is to present the theoretical foundations of the binomial model, implement it and the compare it to other option pricing models.

We will begin by establishing the basic terminology and presenting key theoretical results from financial mathematics.
Once the model's algorithmic implementation is complete, we will explore its capabilities by pricing complex options.
We will analyze how the model's discrete structure allows for the valuation of products that cannot be handled analytically by the Black-Scholes model, such as American options which can be exercised early.
The study will be complemented by the calculation of sensitivity indicators known as \textit{Greeks}, which are essential risk management tools for measuring a portfolio's sensitivity to market volatility.

Finally, we will compare our results with those obtained from the Black-Scholes and Monte Carlo models. 
This comparative analysis will enable us to highlight the binomial model's convergence rate and identiy its limitations when applied to real-world date, thereby preparing for potential extensions involving more sophisticated tree structures.

\subsection{Vocabulary}

Before proceeding with the study of quantitative finance and the algorithms behind the Cox-Ross-Rubinstein model, it is essential to establish a basic vocabulary necessary for a proper understanding of the remainder of the document. 
If you have a financial background, you can skip this part as it will only present the basic key ideas in market finance.


\underline{\textbf{Derivative :}} A derivative is a type of financial contract whose value depends on an asset known as \textbf{underlying asset} and whose price fluctuates in line with changes in the rate or price of that underlying asset.
There are three main categories of derivatives : \textbf{firm commitment products}, \textbf{swaps} and \textbf{contingent products or options} \cite{Derivatives} \cite{ProduitsDerives} \cite{ProduitDerive} 

\underline{\textbf{Forward product :}} This group of derivative products generally includes \textbf{forward} and \textbf{futures}. 
A \textbf{forward} contract is a financial agreement between two entities, agreeing on a defined sale and purchase price at a fixed deadline.
Both entities must honor their parts of the contract even if the situation at maturity is unfavorable to them.
There are some differences between \textbf{forwards} and \textbf{futures} : forward contracts are negotiated over the counter, non-standardized and not very negotiable, whereas futures contracts are negotiated on a stock exchange, are standardized and negotiable. \cite{Derivatives} \cite{ProduitsDerives} \cite{Forward} \cite{Swap} 

\underline{\textbf{Swap :}} Exchange of financial flows between two entities (e.g. floating-rate interest for fixed-rate interest). \cite{Swap}

\underline{\textbf{Option :}} Establishes a right between a buyer and a seller to buy (\textbf{Call}) or sell (\textbf{Put}) an underlying asset at a predetermined price (\textbf{Strike}) by a set expiration date.
This derivative represents a right and is not an obligation.
The \textbf{Call/Put} symmetry allows one to anticipate either a favorable movement in the underlying asset's price (Call) or an unfavorable one (Put).
There are numerous types of options, notably \textbf{European} options that are exercisable only at expiration, and \text{American} options that are exercisable at any time prior to expiration.
These represent the simplest forms, known as \text{Vanilla options}. \cite{VanillaOptions} \cite{Option}
The \textbf{Payoff} refers to the outcome at expiration.

An option is :
\begin{itemize}
    \item \textbf{In the money / ITM} when it strike price is lower than the price of its underlying asset for a call, or higher than the price of its underlying asset for a put.
    \item \textbf{Out of the money / OTM} when its strike price is higher than the price of its underlying asset for a call, or lower than the price of its underlying asset for a put.
    \item \textbf{At the money / ATM} when its strike is equal to the current price of the underlying asset.
\end{itemize}
We will use the following notations :
\begin{itemize}
    \item K : Strike
    \item S : Price of the underlying asset
    \item p : Option Premium
    \item R : Profit at maturity
    \item P : Payoff
\end{itemize}

\bigskip

\begin{definitionbox}[frametitle = {Payoff Formulas}]
\begin{center}
    $P_{Call} = \max(S-K,0)$ \\
    $P_{Put} = \max(K-S,0)$
\end{center}

\end{definitionbox}

We can represent profits at maturity : \vspace{0.025\textheight}
\begin{figure}[htbp]
  \centering
  \begin{subfigure}[b]{0.48\textwidth}
    \centering
    \begin{tikzpicture}
    \begin{axis}[
        width=\textwidth,
        height=0.8\textwidth,
        axis lines = middle,
        xlabel = {$S$},
        ylabel = {$R = \begin{cases} R_{\text{\textcolor{red}{Call - Buyer}}} = \max(S-K, 0) - p \\ R_{\text{\textcolor{blue}{Call - Seller}}} = -\max(S-K, 0) + p \end{cases}$}, 
        ylabel style={at={(axis description cs:0,1.05)}, anchor=south west},
        ymin = -3.5, ymax = 3.5,
        xmin = 0, xmax = 4.5,
        xtick = {2, 2.8},
        xticklabels = {$K$, $K+p$},
        ytick = {-0.8,0.8},
        yticklabels = {$-p$,$p$},
        grid = none
    ]
    \addplot [domain=0:2, samples=2, color=red, thick] {-0.8};
    \addplot [domain=2:4.5, samples=2, color=red, thick] {x - 2.8}; 
    \addplot [domain=0:2, samples=2, color=blue, thick] {0.8};
    \addplot [domain=2:4.5, samples=2, color=blue, thick] {2.8 - x}; 
    \end{axis}
    \end{tikzpicture}
    \caption{Option - Call}
  \end{subfigure}
  \hfill
  \begin{subfigure}[b]{0.48\textwidth}
    \centering
    \begin{tikzpicture}
    \begin{axis}[
        width=\textwidth,
        height=0.8\textwidth,
        axis lines = middle,
        xlabel = {$S$},
        ylabel = {$R = \begin{cases} R_{\text{\textcolor{red}{Put - Buyer}}} = \max(K-S, 0) - p \\ R_{\text{\textcolor{blue}{Put - Seller}}} = -\max(K-S, 0) + p \end{cases}$}, 
        ylabel style={at={(axis description cs:0,1.05)}, anchor=south west},
        ymin = -3.5, ymax = 3.5,
        xmin = 0, xmax = 4.5,
        xtick = {1.2, 2},
        xticklabels = {$K-p$, $K$},
        ytick = {-0.8,0.8},
        yticklabels = {$-p$,$p$},
        grid = none
    ]
    \addplot [domain=0:2, samples=2, color=red, thick] {1.2 - x};
    \addplot [domain=2:4.5, samples=2, color=red, thick] {-0.8}; 
    \addplot [domain=0:2, samples=2, color=blue, thick] {x - 1.2};
    \addplot [domain=2:4.5, samples=2, color=blue, thick] {0.8}; 
    \end{axis}
    \end{tikzpicture}
    \caption{Option - Put}
  \end{subfigure}
  \caption{Profit of Call and Put options}
\end{figure}

\newpage

Options are a key concept to understand qunatitative finance.
It is useful to work through a numerical example to see how it works : 

Consider an industrial buyer of wheat, \textbf{B}, who knows on \textbf{January 1st} that they will need to purchase $n = 1\:000$ tonnes of wheat on \textbf{June 30th}.
On January 1st, the price of wheat is 200\$ per tonne. 
B anticipates that the price of wheat will rise. 
Since B would begin to lose money if the price exceeds 250\$, they decide to purchase $1\:000$ call options with a strike price of 250\$, a premium of $p=10$\$ per tonne and an expiration date of June 30th.

Several scenarios are possible for June 30th : 
\begin{itemize}
    \item \textbf{Case 1}: Wheat trades at 150\$ per tonne.
    The call option is \textbf{OTM}.
    B's prediction did not materialize.
    The call option is now worthless, leading B to let it expire.
    B thus loses $n \cdot p = 1000 \cdot10 = 10\:000$\$.
    B therefore purchases the wheat on the market at 150\$ per tonne, resulting in a total cost of $150 + p = 160$\$ per tonne.

    \item \textbf{Case 2}: Wheat trades at 225\$ per tonne.
    The call option is \textbf{OTM}.
    B's prediction has partly come true but not enough to make the call options worth using.
    The call option is still worthless, leading B to let it expire.
    B thus loses $n \cdot p = 1000 \cdot10 = 10\:000$\$.
    B therefore purchases the wheat on the market at 225\$ per tonne, resulting in a total cost of $225+p=235$\$ per tonne.

    \item \textbf{Case 3}: Wheat trades at 250\$ per tonne.
    The call option is \textbf{ATM}.
    B's prediction has partly come true but not enough to make the call options worth using because the price is equal to the strike.
    Using the call option does not change anything.
    B thus loses $n \cdot p = 1000 \cdot10 = 10\:000$\$.
    B therefore purchases the wheat on the market at 250\$ per tonne, resulting in a total cost of $250 + p = 260$\$ per tonne.
    
    \item \textbf{Case 4}: Wheat trades at 300\$ per tonne.
    The call option is \textbf{ITM}.
    B's prediction came true.
    B uses its call options and spend $250$\$ per wheat tonne.
    Thus, B will have spent $250 + p = 260$\$ per tonne instead of 300\$.
    B saved $(300-(250+p))\cdot n = 40 \cdot 1000 = 40\:000$\$. \\   
  

\end{itemize}
Therefore, the goal of market finance is to optimize strategies in order to generate the highest possible profit while taking risk into account. 
Thus, it is useful to develop mathematical tools that allow for the modeling of market dynamics.

\underline{\textbf{Positions}} : There are two positions when buying or selling options : \\
\begin{itemize}
    \item \textbf{Long Position}
    \begin{description}
        \item[Action] Buying the asset in anticipation of a price increase
        \item[Profit] Making money by reselling at a higher price than the purchase price
        \item[Risk] Risk limited to the amount invested
    \end{description}
    \item \textbf{Short Position}
    \begin{description}
        \item[Action] Borrowing an asset from someone, selling it immediately, waiting for the price to drop and buying back the asset at the lower price to return it
        \item[Profit] Gaining the difference
        \item[Risk] Theoretically infinite risk if the share price rises without limit
    \end{description}
\end{itemize}
